% GENERATED FILE. Edit canonical lessons, then run npm run generate:books.
% canonical-source-hash: fnv1a64:5f0efb063a66f2de
% canonical-lessons: TA-C186-mamanar, TA-C186-katalan, TA-C186-katali, TA-C186-katal, TA-C186-natpu

\chapter{Father-in-law, Boyfriend, Girlfriend, Romantic love, Friendship}
\label{ch:ta-father-in-law-boyfriend-girlfriend-romantic-love-friendship}
\hlchaptermodality{186}

\begin{chapteropening}
\textbf{By the end of this chapter:} \emph{I can say a father-in-law, a boyfriend, a girlfriend, romantic love and friendship.}

Complete the last lesson of chapter 186: I can say a father-in-law, a boyfriend, a girlfriend, romantic love and friendship.
\end{chapteropening}

\section[\texorpdfstring{\ta{மாமனார்} (māmaṉār)}{மாமனார் (māmaṉār)}]{\ta{மாமனார்} (māmaṉār) — a father-in-law}

\label{lesson:TA-C186-mamanar}



\begin{quote}
Before the new one: say the Tamil for news; a message, then the Tamil for an email.
\end{quote}

\subsection*{You'll want to know: \ta{மாமனார்}}
\textbf{\ta{மாமனார்}} — \emph{māmaṉār} — \textquotedblleft{}a father-in-law\textquotedblright{}.

\begin{cousinweb}
\textbf{\ta{மாமனார்}} (\emph{māmaṉār}) is native Dravidian. DEDR 4813 lists Telugu \emph{māma}, \textquotedblleft{}father-in-law\textquotedblright{}. The entry glosses \emph{māmaṉ} \textquotedblleft{}maternal uncle, father-in-law\textquotedblright{}; \emph{-ār} is respectful.
\end{cousinweb}

\begin{grammarlens}[title={What you've built}]
One more word for this chapter.
\end{grammarlens}

\subsection*{Your turn}
\begin{itemize}
\raggedright
  \item \emph{Say it:} \emph{māmaṉār}
  \item \emph{Say it:} \emph{māmaṉār}, once more
  \item \emph{Say it:} say \emph{māmaṉār}
  \item \emph{Recall:} say the Tamil for news; a message, then the Tamil for an email, then say \emph{māmaṉār} again
\end{itemize}

\begin{tcolorbox}[breakable,colback=teal!4,colframe=teal!35!black,title={Before you move on}]
What does \ta{மாமனார்} mean? (\textquotedblleft{}A father-in-law\textquotedblright{}.) Say it once more.
\end{tcolorbox}
\section[\texorpdfstring{\ta{காதலன்} (kātalaṉ)}{காதலன் (kātalaṉ)}]{\ta{காதலன்} (kātalaṉ) — a boyfriend, a (male) lover}

\label{lesson:TA-C186-katalan}



\begin{quote}
Before the new one: say the Tamil for an email, then the Tamil for a father-in-law.
\end{quote}

\subsection*{You'll want to know: \ta{காதலன்}}
\textbf{\ta{காதலன்}} — \emph{kātalaṉ} — \textquotedblleft{}a boyfriend, a (male) lover\textquotedblright{}.

\begin{cousinweb}
\textbf{\ta{காதலன்}} (\emph{kātalaṉ}) is native Dravidian. DEDR 1445 lists Kannada \emph{kādala}, \textquotedblleft{}male friend\textquotedblright{}. It is \emph{kātal}, \textquotedblleft{}love\textquotedblright{}, with \emph{-aṉ} for a man.
\end{cousinweb}

\begin{grammarlens}[title={What you've built}]
One more word for this chapter.
\end{grammarlens}

\subsection*{Your turn}
\begin{itemize}
\raggedright
  \item \emph{Say it:} \emph{kātalaṉ}
  \item \emph{Say it:} \emph{kātalaṉ}, once more
  \item \emph{Say it:} say \emph{kātalaṉ}
  \item \emph{Recall:} say the Tamil for an email, then the Tamil for a father-in-law, then say \emph{kātalaṉ} again
\end{itemize}

\begin{tcolorbox}[breakable,colback=teal!4,colframe=teal!35!black,title={Before you move on}]
What does \ta{காதலன்} mean? (\textquotedblleft{}A boyfriend, a (male) lover\textquotedblright{}.) Say it once more.
\end{tcolorbox}
\section[\texorpdfstring{\ta{காதலி} (kātali)}{காதலி (kātali)}]{\ta{காதலி} (kātali) — a girlfriend, a (female) lover}

\label{lesson:TA-C186-katali}



\begin{quote}
Before the new one: say the Tamil for a father-in-law, then the Tamil for a boyfriend.
\end{quote}

\subsection*{You'll want to know: \ta{காதலி}}
\textbf{\ta{காதலி}} — \emph{kātali} — \textquotedblleft{}a girlfriend, a (female) lover\textquotedblright{}.

\begin{cousinweb}
\textbf{\ta{காதலி}} (\emph{kātali}) is native Dravidian (DEDR 1445). It is \emph{kātal}, \textquotedblleft{}love\textquotedblright{}, with \emph{-i} for a woman.
\end{cousinweb}

\begin{grammarlens}[title={What you've built}]
One more word for this chapter.
\end{grammarlens}

\subsection*{Your turn}
\begin{itemize}
\raggedright
  \item \emph{Say it:} \emph{kātali}
  \item \emph{Say it:} \emph{kātali}, once more
  \item \emph{Say it:} say \emph{kātali}
  \item \emph{Recall:} say the Tamil for a father-in-law, then the Tamil for a boyfriend, then say \emph{kātali} again
\end{itemize}

\begin{tcolorbox}[breakable,colback=teal!4,colframe=teal!35!black,title={Before you move on}]
What does \ta{காதலி} mean? (\textquotedblleft{}A girlfriend, a (female) lover\textquotedblright{}.) Say it once more.
\end{tcolorbox}
\section[\texorpdfstring{\ta{காதல்} (kātal)}{காதல் (kātal)}]{\ta{காதல்} (kātal) — romantic love}

\label{lesson:TA-C186-katal}



\begin{quote}
Before the new one: say the Tamil for a boyfriend, then the Tamil for a girlfriend.
\end{quote}

\subsection*{You'll want to know: \ta{காதல்}}
\textbf{\ta{காதல்}} — \emph{kātal} — \textquotedblleft{}romantic love\textquotedblright{}.

\begin{cousinweb}
\textbf{\ta{காதல்}} (\emph{kātal}) is native Dravidian. DEDR 1445 lists Malayalam \emph{kātal}, \textquotedblleft{}attachment\textquotedblright{} and Kannada \emph{kādal}, \textquotedblleft{}affection, love\textquotedblright{}.
\end{cousinweb}

\begin{grammarlens}[title={What you've built}]
One more word for this chapter.
\end{grammarlens}

\subsection*{Your turn}
\begin{itemize}
\raggedright
  \item \emph{Say it:} \emph{kātal}
  \item \emph{Say it:} \emph{kātal}, once more
  \item \emph{Say it:} say \emph{kātal}
  \item \emph{Recall:} say the Tamil for a boyfriend, then the Tamil for a girlfriend, then say \emph{kātal} again
\end{itemize}

\begin{tcolorbox}[breakable,colback=teal!4,colframe=teal!35!black,title={Before you move on}]
What does \ta{காதல்} mean? (\textquotedblleft{}Romantic love\textquotedblright{}.) Say it once more.
\end{tcolorbox}
\section[\texorpdfstring{\ta{நட்பு} (naṭpu)}{நட்பு (naṭpu)}]{\ta{நட்பு} (naṭpu) — friendship}

\label{lesson:TA-C186-natpu}



\begin{quote}
Before the new one: say the Tamil for a girlfriend, then the Tamil for romantic love.
\end{quote}

\subsection*{You'll want to know: \ta{நட்பு}}
\textbf{\ta{நட்பு}} — \emph{naṭpu} — \textquotedblleft{}friendship\textquotedblright{}.

\begin{cousinweb}
\textbf{\ta{நட்பு}} (\emph{naṭpu}) is native Dravidian. DEDR 3588 lists Kannada \emph{naṇṭu}, \textquotedblleft{}friendship\textquotedblright{} and Telugu \emph{naṇṭu}. It is the same family as \emph{naṇbar}, \textquotedblleft{}friend\textquotedblright{}.
\end{cousinweb}

\begin{grammarlens}[title={What you've built}]
One more word for this chapter.
\end{grammarlens}

\subsection*{Your turn}
\begin{itemize}
\raggedright
  \item \emph{Say it:} \emph{naṭpu}
  \item \emph{Say it:} \emph{naṭpu}, once more
  \item \emph{Say it:} say \emph{naṭpu}
  \item \emph{Recall:} say the Tamil for a girlfriend, then the Tamil for romantic love, then say \emph{naṭpu} again
\end{itemize}

\begin{tcolorbox}[breakable,colback=teal!4,colframe=teal!35!black,title={Before you move on}]
What does \ta{நட்பு} mean? (\textquotedblleft{}Friendship\textquotedblright{}.) Say it once more.
\end{tcolorbox}
